% Authoritative exercise source. Preserve all derivations and diagrams.

\ExerciseStart{G4.02}{Recuperare il periodo binario di Simon}{Guida esplicita}

\BookHeading{Problema, dati e obiettivo}

Un oracolo di Simon associa le coppie \(000/101\), \(001/100\), \(010/111\), \(011/110\) agli esiti rispettivi \(000,001,010,011\). Ricostruire il periodo binario a partire dalle letture 010 e 111. Analizzare un ramo con bersaglio osservato 010.

\ExTheory{\ExRef{c4-simon}{4.03}}{La promessa è \(f(x)=f(x\oplus s)\), con un unico \(s\ne0\). Dopo la misura del bersaglio si ottiene una coppia coerente \((\ket x+\ket{x\oplus s})/\sqrt2\). La Hadamard produce soltanto y tali che \(y\cdot s=0\) in \(\mathbb F_2\).}

\ExPrediction{Le coppie dell’oracolo sono traslazioni binarie. La misura individua un sottospazio di equazioni, non direttamente il periodo. Le equazioni indipendenti ne ridurranno lo spazio dei candidati.}

\ExFigure{\begin{tikzpicture}[>=Stealth,exlabel]
\node[exbox] (n0) at (0.0,0.0) {$000\leftrightarrow101$};\node[exbox] (n1) at (3.6,0.0) {$001\leftrightarrow100$};\draw[exother] (n0)--(n1);\node[exbox] (n2) at (7.2,0.0) {$010\leftrightarrow111$};\draw[exother] (n1)--(n2);\node[exbox] (n3) at (0.0,-1.6) {$011\leftrightarrow110$};\draw[exother] (n2.south)--++(0,-.4)-|(n3.north);
\end{tikzpicture}}{Le quattro coppie condividono la traslazione XOR 101.}


\begin{ExTasks}

\item Verificare la traslazione comune delle quattro coppie.

\item Normalizzare il ramo con bersaglio 010.

\item Calcolare l’ampiezza generica dopo Hadamard.

\item Elencare gli esiti possibili e i loro segni.

\item Risolvere il sistema delle due letture e discutere letture nulle o ripetute.

\end{ExTasks}

\ExSolution

\ExStep{Verificare la promessa}

Le quattro somme XOR sono \(000\oplus101=101\), \(001\oplus100=101\), \(010\oplus111=101\), \(011\oplus110=101\). Ogni valore dell’oracolo compare esattamente due volte.

\ExStep{Condizionare il registro}

La preparazione uniforme contiene otto rami di probabilità \(1/8\). Il bersaglio 010 corrisponde a due rami, quindi ha probabilità \(2/8=1/4\). Dividendo le ampiezze \(1/\sqrt8\) per \(\sqrt{1/4}=1/2\), il primo registro diventa \((\ket{010}+\ket{111})/\sqrt2\).

\ExStep{Ricavare la cancellazione}

\ExEq{A(y)&=\frac14\bigl[(-1)^{010\cdot y}+(-1)^{111\cdot y}\bigr]\\
&=\frac{(-1)^{010\cdot y}}4\bigl[1+(-1)^{101\cdot y}\bigr].}
Se \(101\cdot y=1\), il termine tra parentesi vale zero. Se vale zero, il termine è 2.

\ExStep{Calcolare il supporto}

La condizione è \(y_0\oplus y_2=0\). Gli esiti sono \(000,010,101,111\), con ampiezze rispettivamente \(1/2,-1/2,1/2,-1/2\). I quattro quadrati valgono \(1/4\) e sommano a uno.

\ExStep{Risolvere le equazioni binarie}

La lettura 010 impone \(s_1=0\). La lettura 111 impone \(s_0\oplus s_1\oplus s_2=0\). Sostituendo \(s_1=0\), si ottiene \(s_2=s_0\). I candidati sono \(000\) e \(101\); la promessa esclude il vettore nullo e seleziona 101.

\ExStep{Valutare l’informazione raccolta}

La lettura 000 produce \(0=0\) e non aumenta il rango. Una ripetizione di 010 ripropone la stessa equazione. Nel sottospazio a quattro elementi, dopo una prima equazione non nulla, due dei quattro esiti possibili sono indipendenti da essa. La probabilità di rango due dopo due letture è \((3/4)(2/4)=3/8\).

\par\noindent\begin{minipage}{\linewidth}
\ExFigure{\begin{tikzpicture}[>=Stealth,exlabel]
\node[exbox] (a) at (0,1) {$s_1=0$};\node[exbox] (b) at (0,-.5) {$s_0+s_1+s_2=0$};\node[exbox] (c) at (4,.25) {$s=(t,0,t)$\\$t\in\mathbb F_2$};\draw[exother] (a)--(c);\draw[exother] (b)--(c);
\end{tikzpicture}}{Intersezione delle equazioni sul campo binario. Il parametro t distingue il vettore nullo dal periodo non nullo.}


\ExInterpret{La geometria qui è quella di un sottospazio vettoriale su \(\mathbb F_2\); i segmenti dello schema rappresentano dipendenze logiche. Le cancellazioni delle ampiezze producono le equazioni che il calcolo classico risolve.}

\ExLab{Nella sezione 4.03 scegliere «Segreto 101: equazioni indipendenti». Confrontare con il caso di equazioni ripetute: il numero di letture e il rango della matrice raccolta sono quantità differenti.}


\end{minipage}
